A medical AI scores 98\% on trustworthiness benchmarks but hallucinates drug interactions (reliability failure), succumbs to adversarial prompts (robustness failure), and exhibits demographic bias (fairness failure)---all at nearly identical rates ($r > 0.99$). Are these independent failures requiring separate fixes, or symptoms of a shared root cause? This tight coupling challenges the dominant paradigm in LLM trustworthiness evaluation, where dimensions like reliability, robustness, and fairness are assessed independently using specialized benchmarks such as TrustfulQA, AdvBench, and BOLD. Current practice implicitly assumes independence by reporting separate scores per dimension, leading practitioners to diagnose and fix failures dimension-by-dimension despite mounting evidence that interventions targeting one dimension often affect others.

The consequences of this assumption are both conceptual and practical. Conceptually, independent evaluation presumes that reliability failures (e.g., hallucinations, false statements) arise from mechanisms distinct from robustness failures (e.g., adversarial brittleness) or fairness failures (e.g., demographic bias amplification). Multi-dimensional frameworks like HELM aggregate scores across 50+ benchmarks but do not test whether failures correlate---a model scoring 65\% on TrustfulQA and 68\% on AdvBench could reflect either independent failure modes or a shared underlying deficiency. Practically, this gap misdirects intervention strategies: practitioners may invest compute resources calibrating a model for reliability (targeting hallucinations) when the root cause---poor training data diversity, for instance---simultaneously drives robustness and fairness failures. Without empirical evidence of correlation structure, the field lacks a principled basis for distinguishing shared root causes from dimension-specific issues.

We address this gap through large-scale correlation analysis across 20 LLMs spanning three model size strata (small $<$1B parameters, medium 1-10B, large $>$10B). By computing pairwise Spearman correlations between TrustfulQA (reliability), AdvBench (robustness), and BOLD (fairness) scores, we test whether trustworthiness failures exhibit the moderate correlations ($r > 0.3$) expected from shared root causes, or stronger coupling suggesting dimensional near-redundancy. Our analysis reveals correlations far exceeding initial expectations: all three benchmark pairs exhibit $r > 0.99$ ($p < 2\times10^{-17}$), with this pattern persisting across model scales (stratified analysis shows $r > 0.98$ within small, medium, and large model groups).

This near-perfect correlation challenges the independent dimensions assumption but admits two competing explanations. First, it may indicate that reliability, robustness, and fairness---as operationalized by current benchmarks---measure the same underlying construct (general model quality or unified trustworthiness) rather than orthogonal failure mechanisms. Alternatively, distinct dimensions may exist but 3 benchmarks provide insufficient resolution to distinguish them---correlations could reflect shared confounds (e.g., all benchmarks sensitive to training data diversity) rather than unified constructs. Correlation alone cannot distinguish these explanations without experimental manipulation. Second, when we applied hierarchical clustering to test whether correlations organize into distinct, interpretable failure modes, the analysis revealed perfect cluster stability (100\% bootstrap consistency across 1000 iterations) but poor separation quality (silhouette score = 0.274 $<$ 0.5 threshold). Clusters are reproducible but not well-separated, indicating that our 3-benchmark design cannot resolve the granularity required for failure mode taxonomy---a methodological limitation with implications for future benchmark design.

Building on these findings, we contribute:

\begin{enumerate}
\item \textbf{First large-scale empirical evidence} that trustworthiness dimensions (reliability, robustness, fairness) exhibit near-perfect correlations ($r > 0.99$, $p < 2\times10^{-17}$) across 20 LLMs, far stronger than the moderate effect sizes ($r > 0.3$) hypothesized from shared root causes. This coupling persists across model scales, suggesting either unified constructs or shared confounds affecting all benchmarks---current data cannot distinguish without broader measurement.

\item \textbf{Methodological lesson} for trustworthiness benchmark design: 3-benchmark evaluation is insufficient to resolve distinct failure mode clusters despite perfect stability (100\% bootstrap consistency), with clustering failing separation quality thresholds (silhouette = 0.274). This reveals a concrete design constraint---clustering $k\geq3$ failure modes requires $n\geq3$ benchmarks but silhouette scores below 0.5 indicate poor separation even when clusters are stable.

\item \textbf{Competing explanations} and \textbf{future work roadmap}: The $r > 0.99$ coupling admits two interpretations---(1) Unified Capability Hypothesis, where all three benchmarks measure the same underlying construct, or (2) Insufficient Resolution Hypothesis, where distinct dimensions exist but 3 benchmarks cannot distinguish them. We propose expanding to 5-10 benchmarks (ToxiGen, BBQ, HellaSwag, GSM8K, HumanEval) to test discriminating predictions: Hypothesis 1 predicts $r > 0.95$ persists across most pairs; Hypothesis 2 predicts mean $r$ drops below 0.85 with some pairs $< 0.7$.
\end{enumerate}

Our work shifts the conversation from ``score each dimension independently'' to ``analyze correlation structure first''---a paradigm change with implications for both evaluation methodology and intervention design. If coupling generalizes beyond our 3-benchmark sample, practitioners may benefit from multi-dimensional interventions targeting shared root causes, though intervention validation (testing whether reliability-targeted training also improves robustness/fairness) remains future work. The observed coupling ($r > 0.99$) is robust within our sample, but interpretation remains uncertain pending broader measurement and causal experiments. The taxonomy validation failed due to sample size constraints, and the path forward requires expanding to 5-10 benchmarks to distinguish competing explanations.
